% Guide: using pdf-to-markdown with academic-pdf-reading
% Compile with: xelatex pdf-to-markdown_with_academic-pdf-reading.tex  (run twice)
\documentclass[11pt,a4paper]{article}

\newcommand\skillname{pdf-to-markdown + academic-pdf-reading}
\newcommand\skillversion{}
\newcommand\skilldate{12 October 2026}

% ---------------------------------------------------------------- layout
\usepackage[margin=2.4cm,top=2.7cm,bottom=2.6cm,headheight=14pt]{geometry}
\usepackage{fontspec}
\setmainfont{Latin Modern Roman}
\setsansfont{Latin Modern Sans}
\setmonofont{DejaVu Sans Mono}[Scale=0.82]
\usepackage{microtype}
\usepackage{xcolor}
\definecolor{accent}{HTML}{1F5F7A}
\definecolor{codebg}{HTML}{F2F5F7}
\definecolor{muted}{HTML}{6B7780}
\usepackage{titlesec}
\titleformat{\section}{\sffamily\Large\bfseries\color{accent}}{\thesection}{0.8em}{}
\titleformat{\subsection}{\sffamily\large\bfseries}{\thesubsection}{0.7em}{}
\titlespacing*{\section}{0pt}{1.5em}{0.5em}
\titlespacing*{\subsection}{0pt}{1.1em}{0.35em}
\usepackage{enumitem}
\setlist{itemsep=0.2em, topsep=0.35em}
\usepackage{booktabs,tabularx,array}
\renewcommand\tabularxcolumn[1]{>{\raggedright\arraybackslash}p{#1}}
\setlength\extrarowheight{1pt}
\usepackage[most]{tcolorbox}
\tcbuselibrary{listings}
\newtcblisting{cmd}{listing only, colback=codebg, colframe=codebg, boxrule=0pt,
  arc=2pt, left=6pt, right=6pt, top=2pt, bottom=2pt, before skip=6pt, after skip=8pt,
  listing options={basicstyle=\ttfamily\small, breaklines=true, columns=fullflexible,
  keepspaces=true, upquote=true, breakindent=1.5em}}
\newtcolorbox{glance}{enhanced, colback=accent!5, colframe=accent, boxrule=0.6pt,
  arc=2pt, left=9pt, right=9pt, top=6pt, bottom=6pt,
  fonttitle=\sffamily\bfseries, title=At a glance}
\newtcolorbox{note}{enhanced, colback=accent!4, colframe=accent, boxrule=0pt,
  leftrule=2.5pt, arc=0pt, left=8pt, right=8pt, top=5pt, bottom=5pt,
  before skip=8pt, after skip=8pt}
\usepackage{fancyhdr}
\pagestyle{fancy}
\fancyhf{}
\fancyhead[L]{\sffamily\small\color{muted}\skillname}
\fancyhead[R]{\sffamily\small\color{muted}Guide}
\fancyfoot[C]{\sffamily\small\color{muted}\thepage}
\renewcommand\headrulewidth{0.3pt}
\usepackage[colorlinks=true, linkcolor=accent, urlcolor=accent]{hyperref}
\hypersetup{pdftitle={Using pdf-to-markdown with academic-pdf-reading},
  pdfsubject={Claude skills guide}}
\DeclareUrlCommand\code{\urlstyle{tt}}
\newcommand\step[1]{\textsf{\textbf{#1}}}

\begin{document}

\thispagestyle{plain}
\begin{flushleft}
{\sffamily\small\color{muted}\MakeUppercase{Claude skills \textbullet{} Guide}}\\[4pt]
{\sffamily\fontsize{22}{27}\selectfont\bfseries\color{accent}Using pdf-to-markdown with\\ academic-pdf-reading}\\[6pt]
{\sffamily\large Accurate equations at lower token cost\quad\textcolor{muted}{|}\quad \skilldate}\\[2pt]
{\color{accent}\rule{\linewidth}{0.8pt}}
\end{flushleft}

\begin{glance}
\textbf{In one sentence:} convert equation-heavy PDFs to Markdown once with
\textbf{pdf-to-markdown}, and \textbf{academic-pdf-reading} will read the transcription as
exact source, so Claude gets image-verified equations without paying for page images in every
session.
\end{glance}

\section{The trade-off}
academic-pdf-reading is accurate because it never takes mathematics from a PDF's text layer.
Every sign, index, bar, prime and typeface is read from the page image: a full page view,
then an upscaled crop of each line that matters. That accuracy has a cost, paid on every
question about an equation and paid again in every new session, because nothing Claude saw in
one conversation carries over to the next.

\section{Transcribe once, read as source}
pdf-to-markdown spends the image-reading effort a single time. Claude transcribes every formula
from page images into LaTeX. Scripts then check the syntax, equation tags and page markers,
and each displayed equation is proofread side by side against its crop from the PDF. The
result is a text-only \code{NAME.md} plus \code{NAME_figures.pdf}.

academic-pdf-reading recognises a transcription by its first-line legend and treats it like
LaTeX source:
\begin{itemize}
\item it quotes equations as written, without re-checking them against page images;
\item it cites pages from the transcription's page markers, including printed page numbers;
\item it opens the figures PDF only for questions about what a figure shows;
\item it escalates only when the transcription is contradicted by a derivation or a companion
  paper, is internally inconsistent, or is structurally defective. Claude never corrects a
  transcription silently.
\end{itemize}

\section{What it costs}
\begin{center}
\small
\begin{tabularx}{\linewidth}{@{}l X X@{}}
\toprule
 & \textbf{Reading the PDF directly} & \textbf{Reading a transcription} \\
\midrule
One-time & none & conversion: about 2{,}500--3{,}000 tokens per page, plus about 1{,}500 per
  proofreading sheet (four equations each) \\
Each read of a page & a page image (up to about 1{,}600 tokens) plus crops of the relevant
  lines, often several thousand tokens in total & the page's Markdown text, typically well
  under 1{,}000 tokens \\
Each new session & the same again & the same text cost; no images \\
\bottomrule
\end{tabularx}
\end{center}
\noindent\textcolor{muted}{\small Per-read figures are approximate and depend on page density
and on how many lines need cropping.}

\medskip
\noindent A page costs roughly as much to convert as one or two careful image readings of it.
Converting therefore pays for itself whenever a document will be consulted for more than a
handful of equations, or across more than one session. After that, every read is cheaper.

\section{When to convert}
\textbf{Convert} equation-heavy papers, lecture notes and textbook chapters that you will keep
returning to, especially documents stored in a Project.

\medskip\noindent\textbf{Leave as PDF} in these cases:
\begin{itemize}
\item a one-off question about a single equation, where direct reading is cheaper;
\item figure-centric documents, since text inside figures is not transcribed;
\item documents whose layout itself is the question.
\end{itemize}
Convert whole books one chapter per conversation.

\section{Workflow}
\begin{enumerate}
\item \step{Convert.} In a chat, attach the PDF and ask Claude to convert it to Markdown.
  For long documents, give a page range.
\item \step{Download} \code{NAME.md} and, if present, \code{NAME_figures.pdf}.
\item \step{Upload} both to the Project's knowledge. The figures PDF can be omitted if no
  figure will need to be seen, because every caption is already in the Markdown.
\item \step{Keep or remove the original PDF.} If both are present, Claude reads the
  transcription for content and opens the PDF only for questions about the page itself.
\item \step{Update the protocol}, if the Project has one. The staleness check reports the
  document as \emph{converted}; accept Claude's offer to rewrite its entry.
\end{enumerate}

\begin{note}
\textbf{Where the trust lies.} A transcription moves verification from reading time to
conversion time. Displayed equations are proofread against the PDF. Inline math is checked by
re-reading the strips with mandatory zooms on accents and indices, not by side-by-side sheets.
Any symbol left ambiguous at conversion is named in the conversion report: keep that report.
\end{note}

\end{document}
